\section{From one block to a finite stream}\label{h:stream:section}

The block estimate applies to the success function obtained by starting
the learner from a specified state. We first construct these fresh
continuations and show that a block has only a bounded number of possible
destinations. This count will let us propagate the local bound
\eqref{h:local:bound} backward through the stream.

For now fix the shared rule seed and suppose that all finite transition
and stopping kernels are Borel and defined at every input pair. Future
rows and fresh random choices are independent of the past. We take the
action at the horizon to be terminal from every state.
Section~\ref{h:stream:reduction} justifies these conventions for the full randomized,
completed-measurable model of Definition~\ref{h:model}, without changing
its prior expected terminal score.

\subsection{Fresh continuations and local destinations}

Write $Q_t$ for the set of persistent states at sample index $t$, with
$|Q_t|\le2^M$. For each terminal state $w\in Q_k$, write
$\mathsf Q_{k,w}$ for its output probability law on $S^{d-1}$;
with the seed fixed, this law depends only on $k,w$.
Fix a Borel terminal score
$\ell:S^{d-1}\times S^{d-1}\to[0,1]$. Starting from a specified
state $u\in Q_t$ before its sample-free stopping check, define
$h_{t,u}(s)$ to be the expected terminal score
$\mathbb E[\ell(s,\widehat s)]$ with signal $s$ and fresh future
samples and randomness. This definition also applies to states the
original run never reaches.

A block from index $t$ to index $t+m$ includes the stopping check at
its entrance. An active exit is cut before the check at $t+m$, which
the continuation performs. A terminal tag produced by the last
transition has already stopped at $t+m$. These phase conventions use
only the state and the index.

\begin{lemma}[The destinations of a block]\label{h:stream:categories}
For a fixed-seed Borel learner with the state sets, output laws, and
phase convention just specified, and the Borel terminal score $\ell$, let
$0\le t<T$, $u\in Q_t$, and $1\le m\le T-t$. There are an integer
\begin{equation}\label{h:stream:destinations}
 1\le N\le(m+2)2^M,
\end{equation}
Borel routing functions
$g_j:\mathbb R^{m\times d}\times\mathbb R^m\to[0,1]$ with
$\sum_jg_j\le1$, and Borel functions
$h_j:S^{d-1}\to[0,1]$ such that, for every $s\in S^{d-1}$,
\begin{equation}\label{h:stream:category-representation}
 h_{t,u}(s)=\sum_{j=1}^N h_j(s)
       \int g_j(A,As)\,d\gamma_m(A).
\end{equation}
Here $\gamma_m$ is the law of an $m$-by-$d$ standard Gaussian matrix.
Each destination is either a terminal state $w\in Q_k$ at an index
$k\in\{t,\ldots,t+m\}$, with score function
\[
 h^{\mathrm{out}}_{k,w}(s)
 =\int_{S^{d-1}}\ell(s,v)\,\mathsf Q_{k,w}(dv),
\]
or a continuing state $w\in Q_{t+m}$ before its stopping check, with
score function $h_{t+m,w}(s)$. The terminal output law
$\mathsf Q_{k,w}$ depends only on its state and index. The value
$g_j(A,b)$ is the probability of destination $j$ when the block data
are $(A,b)$.
\end{lemma}

\begin{proof}
The everywhere-defined kernels and the forced terminal action at $T$
give a finite scored experiment from every specified state. Its score
function is Borel in the signal. Indeed, the map
$(s,x)\mapsto(x,\langle x,s\rangle)$ is continuous, and each path
probability is a finite product of Borel coordinate probabilities.
Summing over the finitely many paths and integrating against the fixed
laws of future rows preserves Borel measurability. The same argument
applied to $\ell(s,v)$ and each fixed output law shows that
$h^{\mathrm{out}}_{k,w}$ is Borel. The integration fact follows first
for indicators of rectangles and then for nonnegative Borel functions
by the monotone class theorem and monotone convergence; see also the
composition of kernels in \cite[Section~2.3]{GagnePanangaden2023}.

Pre-generate all $m$ rows of the block, including those after an early
stop. At each $k=t,\ldots,t+m-1$, perform the sample-free stopping
check and, if still active, process the next pair through the transition
to $k+1$. A terminal tag from that transition stops at $k+1$.
An active exit from the final transition is passed to the continuation
before the stopping check at $t+m$, so that check is performed exactly
once. If $t+m=T$, the continuation performs the forced terminal action.

A terminal destination records its local stopping offset
$i\in\{0,\ldots,m\}$ and its terminal state in $Q_{t+i}$.
The offset $i=0$ includes an immediate stop, and $i=m$ can arise
from the last transition. A transition tag and a sample-free stop at
the same state and index have the same output law, so their outcomes
can be merged. Continuing destinations form one additional copy of
$Q_{t+m}$, distinguished from the terminal destinations at that index.
Consequently the number of destinations is at most
\[
 \sum_{i=0}^m|Q_{t+i}|+|Q_{t+m}|\le(m+2)2^M.
\]
The block's entrance index is fixed, so local offsets determine all
absolute stopping indices; the count has no factor depending on $T$.

For arbitrary data $(A,b)$, each destination probability $g_j(A,b)$
is a finite sum over paths of products of Borel coordinate
probabilities. It is therefore Borel. The probabilities are
nonnegative and sum to one when all destinations are included.
A path that stops early is simply independent of the unused rows.
After a continuing destination is fixed, independent future rows and
fresh choices give its prescribed continuation score. After a terminal
destination is fixed, the prescribed output law gives its terminal
score. Conditioning first on the block data and then on its destination
proves \eqref{h:stream:category-representation}. Because the kernels
are defined everywhere, this equality holds for every signal and every
starting state, rather than only at states reached by the original run.
\end{proof}

The angular and chordal success scores are, respectively,
\[
 \ell(s,v)=\mathbf 1_{\{\arccos\langle s,v\rangle\le\epsilon\}},
 \qquad
 \ell(s,v)=\mathbf 1_{\{\|s-v\|\le\epsilon\}}.
\]
Both are Borel. The backward induction below uses the angular score.

\subsection{Borel versions and the original randomized experiment}
\label{h:stream:reduction}

The fresh continuations require kernels defined even at inputs and
states absent from the original run. The model permits transition
coordinates measurable in a completed sigma field. Under a dominated
fresh-pair law, these coordinates admit everywhere-defined Borel
versions that preserve the prior experiment. The following statement
keeps the prior general for the later cube application.


\begin{lemma}[Borel versions under a dominated sample law]
\label{h:stream:borel}
Let $d\ge2$, let $\pi$ be a Borel probability measure on $S^{d-1}$,
and let $X\sim\gamma_1=N(0,I_d)$ be independent of $S\sim\pi$.
Suppose that
\begin{equation}\label{h:stream:pair-domination}
 \mathcal L(X,\langle X,S\rangle)\ll
 \lambda,\qquad \lambda(dx,dy)=\gamma_1(dx)\,dy
\end{equation}
on $\mathbb R^d\times\mathbb R$.
Fix an integer horizon $T\ge0$ and finite state sets at every sample
index. For each state and index, suppose every coordinate of the
finite vector of transition and stopping probabilities is measurable
in the $\lambda$-completion of the Borel sigma field. A terminal output
law is a Borel probability measure on $S^{d-1}$ indexed only by its
terminal state and stopping index. Decisions made without a new sample
also depend only on the state and index.

There are everywhere-defined Borel transition and stopping kernels,
with the same state sets and terminal output laws, whose $\pi$-prior
joint law of the signal, state path, stopping index and output is the
same as for the original rules. The domination
\eqref{h:stream:pair-domination} holds when $\pi=\sigma$ is uniform
spherical probability, for every $d\ge2$.
\end{lemma}

\begin{proof}
We first construct versions for a fixed set of rules. If
$f:\mathbb R^d\times\mathbb R\to[0,1]$ is measurable in the
$\lambda$-completion, each dyadic superlevel set
$\{f\ge k2^{-n}\}$, $1\le k\le2^n$, agrees with a Borel set outside
a Borel $\lambda$-null set, by the description of completion in
Tao~\cite[Exercise~1.4.26]{TaoMeasure2011}. Replace these superlevel sets by their
Borel representatives in
\[
 f_n=2^{-n}\sum_{k=1}^{2^n}\mathbf 1_{\{f\ge k2^{-n}\}}.
\]
The resulting Borel functions have a Borel $\limsup$ equal to $f$
outside the union of the countably many null sets, since $f_n\to f$.
Thus each coordinate of each finite kernel has a Borel version.
For one state and index, choose all of its coordinate versions
simultaneously. The Borel set on which the resulting vector is not
nonnegative with sum one is null; on this set replace the vector by a
fixed point mass. The result is a Borel probability vector at every
sample pair. There are finitely many coordinates, states and indices,
so one Borel $\lambda$-null set $Z$ contains every pair where any
coordinate was changed. Sample-free decisions and terminal output
laws require no such modification.

Pre-generate all $T$ input rows, including rows after an early stop.
Every pair $(X_t,\langle X_t,S\rangle)$ has the marginal law in
\eqref{h:stream:pair-domination}, so
\[
 \mathbb P\{(X_t,\langle X_t,S\rangle)\in Z\}=0.
\]
For a state $u$ chosen from earlier observations, the event that the
learner is in $u$ before sample $t$ and this pair lies in $Z$ is a
subset of the displayed null event. We do not need a density for a
sample conditioned on a particular state. A finite union over $t$
shows that no changed pair is visited almost surely. Couple the
original and modified finite choices by the same uniform random
numbers. Starting from the same initial state, induction makes their
states and stopping decisions equal almost surely. At an equal
terminal state and index the output laws are identical, so use their
diagonal coupling. This proves the asserted joint-law equality. When
$T=0$ there is no sample pair to modify and the same conclusion is
immediate.

It remains to check the spherical prior. Conditional on a nonzero
$X=x$, rotate coordinates so that $x/\|x\|$ is the first coordinate
vector. Write a uniform sphere point as $G/\|G\|$ for a standard
Gaussian vector $G=(G_1,G_\perp)$. When $d\ge2$,
$R=\|G_\perp\|>0$ almost surely and is independent of $G_1$.
For each fixed $R>0$, the map
\[
 g\longmapsto \frac{g}{\sqrt{g^2+R^2}}
\]
is a smooth strictly increasing bijection from $\mathbb R$ to
$(-1,1)$ with nonzero derivative. The image of the Gaussian law of
$G_1$ therefore has a Lebesgue density. Mixing over $R$ and then
multiplying by $\|x\|$ shows that $\langle x,S\rangle$ is absolutely
continuous for every $x\ne0$. The point $x=0$ has $\gamma_1$-measure
zero. Fubini's theorem applied to the sections of a
$\lambda$-null Borel set now proves
\eqref{h:stream:pair-domination}.
\end{proof}

The equality in Lemma~\ref{h:stream:borel} concerns the experiment under
the chosen prior. It need not preserve behavior at every exceptional
individual signal. In particular, an every-signal success premise will
be averaged over the prior before versions are chosen.

We now return to the model's original jointly measurable randomized
experiment. Fix any Borel prior $\pi$ satisfying
\eqref{h:stream:pair-domination}, and draw the signal from $\pi$
independently of the rows and the model's data-independent randomness.
The preceding lemma supplies the domination condition for the
uniform-sphere prior. Write $W$ for the shared rule seed, taking $W$ constant
when no seed is used, and $I$ for its initial state. The observable
signal, rows, state path, stopping index, and output take values in
standard Borel spaces. For a Borel terminal score $\ell$, its original
conditional expectation is a measurable function of $W$ and integrates
to the mixed expected score. Since $(I,W)$ is jointly independent of the signal and all
pre-generated rows, conditional on almost every $W=w$ those inputs
retain their joint law and $I$ is independent of them. The dependence
allowed between $I$ and $W$ does not affect this conclusion.

Fix such a seed value for which the model's per-seed kernel hypotheses
also hold. Lemma~\ref{h:stream:borel}
replaces its completed transition and stopping coordinates by
everywhere-defined Borel kernels while preserving the original
conditional $\pi$-prior law. Keep its terminal output laws. At index $T$,
force the sample-free action to stop at every state, retaining each
prescribed state/index output law. Assign a fixed unit output only where
no law was specified, a case absent from original terminating paths.
This does not change the original conditional
experiment, which already stops by $T$. It makes the horizon convention
in the continuation construction valid even at states absent from that
experiment.

The preceding continuation lemma now applies to this fixed-seed Borel
learner with any Borel terminal score $\ell$. A numerical score bound
uniform over its initial states can be averaged against the conditional
law of $I$. If the bound is also uniform in the fixed seed, the final
seed average uses the original conditional expected scores. Borel
versions may be chosen separately for each seed; no jointly measurable
choice of these replacements is required.

\subsection{The quantitative success bound}

We apply the block bound to the fresh continuations just constructed.
At every step, all terminal destinations have the original terminal
local bound, and all continuing destinations have the inductive bound
at that boundary. The following statement records the resulting
finite-horizon estimate before we extract the sample lower bound.

\begin{proposition}[Success after a bounded number of samples]
\label{h:stream:success}
There are absolute constants $C,d_0>0$ with the following property.
For $d\ge d_0$, every learner in the stated finite-state model with
$M\ge0$ bits and an integer deterministic horizon $T\ge0$ has, for
$S\sim\sigma$ independent of its fresh Gaussian rows,
\begin{equation}\label{h:stream:success-bound}
 \mathbb P\{\arccos\langle\widehat s,S\rangle\le\epsilon\}
 \le
 \min\left\{1,
 \left[2\epsilon\exp\!\left(
 C\left(1+\frac{M}{d^2}\right)
 \left\lceil\frac{T}{q}\right\rceil\right)\right]^a
 \right\},
\end{equation}
where $0<\epsilon\le1/10$ and
$a=(d-1)/2$, $q=\lfloor(d-1)/8\rfloor$ as in
\eqref{h:local:parameters}. The probability includes data-independent
initialization, the shared seed, and fresh random choices.
\end{proposition}

\begin{proof}
Use the angular score in Lemma~\ref{h:stream:categories}. Fix a seed
and a prior-preserving Borel replacement as in
Section~\ref{h:stream:reduction}, with prior $\sigma$.
If $T=0$, its output is a signal-independent mixture of terminal laws.
Lemma~\ref{h:local:caps} gives local mass at scale $\epsilon$.
The radius-two ball about a sphere point covers $S^{d-1}$, so total
success is at most $(2\epsilon)^a$, as required when
$\lceil T/q\rceil=0$. We may now suppose $T\ge1$.

Increase the absolute threshold $d_0$ so that
$q\ge2$, $q\ge d/16$, and $a\ge d/3$. Partition the $T$ sample
indices into $B=\lceil T/q\rceil$ consecutive nonempty blocks, each
of length at most $q$. At the final boundary every continuation is
a terminal output function and obeys \eqref{h:local:bound} at scale
$\epsilon$. Suppose that at the end of one block, all continuing
functions obey that bound at a common scale $R\ge\epsilon$. Its
terminal destinations also obey the bound at scale $R$, by
Lemma~\ref{h:local:caps} and monotonicity in the scale.

For a block of length $m\le q$, Lemma~\ref{h:stream:categories}
places every preceding continuation in the form required by
Lemma~\ref{h:block:bound}, with
$N\le(m+2)2^M\le(q+2)2^M$. That block lemma multiplies the scale
by at most
\[
 \exp(C_0d/a)N^{1/(aq)}
\]
for an absolute $C_0$. Its logarithm is at most
\begin{equation}\label{h:stream:multiplier}
 \frac{C_0d}{a}+
 \frac{M\log2+\log(q+2)}{aq}
 \le C\left(1+\frac{M}{d^2}\right)=:\Lambda,
\end{equation}
where $d/a\le3$, $1/(aq)\le48/d^2$, and
$\log(q+2)/(aq)\le48\log(d+2)/d^2$ allow an absolute $C$.
Backward induction gives \eqref{h:local:bound} at scale
$R=\epsilon e^{\Lambda B}$ for every initial state. Averaging over
the conditional initial-state law preserves this bound, because that
law is independent of the signal and rows. Taking the radius-two ball
gives total success at most $(2\epsilon e^{\Lambda B})^a$, and a
probability is also at most one.

By the reduction in Section~\ref{h:stream:reduction}, this is a
uniform bound on the original conditional success probabilities for
almost every seed. Averaging them proves
\eqref{h:stream:success-bound}; no joint choice of the Borel
replacements is used.
\end{proof}
